\documentclass{article}
\usepackage[margin=28mm]{geometry}
\usepackage{amsmath,amssymb}
\usepackage{xurl}
\usepackage[hidelinks]{hyperref}
% Shared commands for notes in docs/paper-gaps/.

\DeclareUnicodeCharacter{2080}{\ifmmode{}_{0}\else\textsubscript{0}\fi}
\DeclareUnicodeCharacter{2081}{\ifmmode{}_{1}\else\textsubscript{1}\fi}
\DeclareUnicodeCharacter{2082}{\ifmmode{}_{2}\else\textsubscript{2}\fi}
\DeclareUnicodeCharacter{2099}{\ifmmode{}_{n}\else\textsubscript{n}\fi}

\newcommand{\Lean}{\textsc{Lean}}
\ExplSyntaxOn
\NewDocumentCommand{\leanid}{m}{
  \ifmmode
    \textsf{\detokenize{#1}}
  \else
    \group_begin:
    \urlstyle{sf}
    \tl_set:Nn \l_tmpa_tl {#1}
    \tl_replace_all:Nnn \l_tmpa_tl {\_} {_}
    \exp_args:NV \path \l_tmpa_tl
    \group_end:
  \fi
}
\ExplSyntaxOff
\providecommand{\tr}{\operatorname{tr}}
\newcommand{\tnleanIssueBase}{https://github.com/LionSR/TNLean/issues/}
\newcommand{\tnleanPrBase}{https://github.com/LionSR/TNLean/pull/}
\newcommand{\tnleanDocBase}{https://sirui-lu.com/TNLean/docs/}
\newcommand{\ghissue}[1]{\href{\tnleanIssueBase#1}{GitHub issue \##1}}
\newcommand{\ghpr}[1]{\href{\tnleanPrBase#1}{PR \##1}}
\newcommand{\doclink}[1]{\href{\tnleanDocBase#1.html}{\path{#1}}}

% Dirac notation and the identity operator, as in blueprint/src/macros/common.tex.
% \providecommand keeps note-local definitions authoritative.
\usepackage{dsfont}
\providecommand{\ket}[1]{|#1\rangle}
\providecommand{\bra}[1]{\langle #1|}
\providecommand{\braket}[2]{\langle #1 | #2 \rangle}
\providecommand{\ketbra}[2]{|#1\rangle\!\langle #2|}
\providecommand{\Id}{\mathds{1}}
\providecommand{\one}{\mathds{1}}
\providecommand{\C}{\mathbb{C}}
\providecommand{\MN}[1]{M_{#1}(\C)}
\providecommand{\MD}{M_D(\C)}

% Verdict marker: \gapnote{<kind>}{<status>}. Kind is the policy
% classification (clarification, local-correction, scope-restriction,
% unfaithful, false-source, open-gap); status is open, wip (elimination
% underway), resolved, or historical. Machine-read by the site index;
% prints a verdict line.
\newcommand{\gapnote}[2]{\par\noindent\textbf{Verdict:} #1 (#2).\par}

\providecommand{\gapnote}[2]{\par\noindent\textbf{Verdict:} #1 (#2).\par}
\title{Geometric Growth in Hastings' Local Approximation and Ground-Space Basis}
\author{}
\date{October 28, 2026}
\begin{document}
\maketitle
\gapnote{scope-restriction}{resolved}

\section{The hypotheses in the source}

Hastings, \emph{Solving gapped Hamiltonians locally}~\cite{Hastings2006Solving},
\href{https://arxiv.org/abs/cond-mat/0508554v3}{arXiv:cond-mat/0508554v3},
source \path{proj3new.tex}, lines 169--190, assumes a metric on the sites,
interaction range $R$, term norms at most $J$, and at most $S$ sites within
distance $R$ of each site. It does not assume polynomial growth of balls or
spheres. A spectral gap and a bounded ground-space dimension do not imply
such geometric growth either.

Two distinct uses of growth must be separated. The earlier proof of the
exponential form of the truncated approximation, eq.~(facld), assumes a
polynomial bound on the omitted shell. A weighted-tail argument now removes
that assumption and proves this conclusion on the source's arbitrary finite
metric spaces. In contrast, the summed basis estimate,
eq.~(oibu), is proved under the uniform sphere bound
\begin{equation}
 \#\{y:d_G(x,y)=r\}\le K_g(r+1)^k.
 \label{eq:hastings05_geometric_growth_scope_spheres}
\end{equation}
Here the constants must be independent of the member of a family of finite
systems. Choosing $K_g$ separately as the total number of sites in each
finite graph does not prove a size-independent estimate.

Polynomial growth is an additional hypothesis. The source's use of the word
``lattice'' does not supply it: it introduces an arbitrary metric at lines 184--187,
and explicitly specializes to a translation-invariant, $d$-dimensional
system when needed at lines 1603--1642. The polynomial-volume interpretation
of the thermal index count is separately discussed in
\path{hastings05_planar_domain_scope.tex}.

\section{A bounded-degree example}

Let $T_L$ be the complete rooted binary tree of depth $L$, with graph
distance. Every radius-one ball has at most four vertices, whereas the
sphere of radius $r\le L$ about the root has $2^r$ vertices. Consequently
no constants $K_g,k$ independent of $L$ satisfy
\eqref{eq:hastings05_geometric_growth_scope_spheres}.

This example is compatible with the other physical hypotheses. Put a qubit
at every vertex and take
\[
 H_L=\sum_{v\ne\mathrm{root}}|1\rangle\langle1|_v,
 \qquad H_{\mathrm{root}}=0.
\]
The terms have norm at most one and range at most one. The ground space has
dimension two and the gap is one, uniformly in $L$. Thus all the source
hypotheses still fail to imply polynomial growth. In this particular
Hamiltonian every ground-space compression of a term is zero. The example
therefore refutes the missing implication about geometry; it is \emph{not}
a counterexample to the source's conclusions.

\section{Why pointwise commutator decay alone is insufficient}

Even in dimension two, an abstract Hermitian family with exponential
commutator decay need not admit a uniformly bounded sum of off-diagonal
norms on these trees. Fix $0<\gamma<\log2$ and $L\ge1$, and split the
$2^L$ leaves into two equal classes. On a
fixed two-dimensional space put
\[
 B_i=e^{-\gamma L}\sigma_x\quad\hbox{on the first class},\qquad
 B_i=e^{-\gamma L}\sigma_z\quad\hbox{on the second class},
\]
and put $B_i=0$ at internal vertices. Since the tree diameter is $2L$,
\[
 \|B_i\|\le1,\qquad
 \|[B_i,B_j]\|\le2e^{-2\gamma L}
       \le2e^{-\gamma d_G(i,j)}.
\]
In a qubit basis with Bloch axis $u$, the off-diagonal norms of $\sigma_x$
and $\sigma_z$ are $\sqrt{1-u_x^2}$ and $\sqrt{1-u_z^2}$.
Their sum is at least one, because the sum of their squares is at least
one. Thus every basis has
\[
 \sum_i o(B_i)\ge2^{L-1}e^{-\gamma L},
\]
which is unbounded when $0<\gamma<\log2$.
This abstract family shows that
pointwise decay, finite ground-space dimension and a local norm budget do
not by themselves remove the growth hypothesis from the present basis
argument. A separate physical construction on a $256$-ary tree now
disproves the unrestricted source claim: it has a two-dimensional ground
space, a uniform positive gap and bounded radius-one interaction terms,
but every basis has $\sum_i o(M_i)\ge2^{L-2}$. The actual localized terms
equal the interaction terms at $l_{\rm proj}=R=1$. Its complete mathematical
proof is in
\href{https://sirui-lu.com/TNLean/paper-gaps/hastings05_tree_basis_counterexample.pdf}
{the tree counterexample note}. That family and its gap proof have not
been formalized in Lean.

\section{Current scope and elimination plan}

The finite-dimensional spectral splitting and recursive construction of an
orthonormal basis are proved without geometric hypotheses. The physical
summed estimates in ground-space dimension two and in arbitrary fixed
dimension use \eqref{eq:hastings05_geometric_growth_scope_spheres}. The
direct estimate for the actual truncated terms also fixes their support
parameters. The perturbative estimate permits growing supports when the
total truncation error stays bounded; its original-term basis estimate
still assumes \eqref{eq:hastings05_geometric_growth_scope_spheres}.
The earlier pairwise observable-decay estimate carries the same growth
hypothesis because it uses partial-trace localization. The direct numerical
finite-support propagation estimate removes that hypothesis for individual
pairs: \path{LowEnergyMetricCommutator.lean} proves eq.~(blecm) for the
actual truncated operators on every finite metric space, with the explicit
support cardinality prescribed in the source. This pointwise conclusion
does not justify a uniform sum over all sites. None of the summed estimates
is a formalization of eq.~(oibu) on every metric lattice of lines 169--190. The counterexample
shows that the extra growth assumption is a correction of a false
unrestricted assertion, not merely a hypothesis awaiting removal.

\subsection{The general metric truncation estimate}

The general metric form of eq.~(facld) is now proved in
\path{GaussianFilteredMetricTruncation.lean}. The constants are chosen before
the finite metric space and Hamiltonian. Here is the argument and one explicit
choice of constants. Let $R>0$, $J\ge0$ and $\Delta>0$. Join two distinct sites
when their physical distance is at most $R$. The resulting graph $G$ has degree
at most $S$, every interaction support has graph diameter at most two, and the
sum of norms of terms meeting a site is at most $SJ$. Set
\[
 \mu=\log(S+1)+1,\qquad
 v=4S^2J e^{4\mu},\qquad
 c_1=\frac{1}{4R(v+1)}.
\]
In particular $c_1$ is independent of the spectral gap. For projection radius
$\ell\ge R$, retain the source's actual set
$\{j:d(i,j)\le\ell-R\}$ and let $m=\lceil(\ell-R)/R\rceil$.
Every omitted anchor has graph distance at least $m$ from $i$.
The weighted omitted-term estimate is
\[
 \sum_{j\text{ omitted}}\|H_j\|e^{-\mu d_G(i,j)}
 \le \frac{J e^{-m/2}}{1-e^{-1/2}}
 \le \frac{J e^{1/2}}{1-e^{-1/2}}e^{-\ell/(2R)}.
\]
A disconnected anchor contributes zero to this sum. This estimate follows by
counting graph balls by $(S+1)^r$ and summing a geometric series. It does not
bound the cardinality of a physical metric ball.

Apply the Lieb--Robinson estimate separately to each omitted term in the
Duhamel formula and choose $T=c_1\ell$, $q=c_1\ell\Delta$. Write
\[
 \gamma=\frac1{2R}-vc_1>0,\qquad
 A=\frac{4(S+1)e^{5\mu}Jc_1e^{1/2}}{1-e^{-1/2}},\qquad
 \delta=\min(c_1\Delta,\gamma),\qquad
 C=3+\frac{2A}{e\delta}.
\]
Gaussian integration and the spectral leakage estimate give
\begin{align*}
 \|H_i\|e^{-q/2}+\|\widetilde H_i^0-M_i\|
 &\le 3\|H_i\|e^{-c_1\Delta\ell/2}
       +A\|H_i\|\ell e^{-\gamma\ell}\\
 &\le \|H_i\|C e^{-\delta\ell/2}.
\end{align*}
Consequently, for every ground vector $\psi$,
\[
 \|P_{\rm high}M_i\psi\|
 \le JC e^{-\delta\ell/2}\|\psi\|.
\]
The actual $M_i$ is Hermitian, is supported in the physical $\ell$-ball and
has norm at most $J$. Their sum satisfies
\[
 \Bigl\|\sum_iM_i-H\Bigr\|\le |\Lambda|JC e^{-\delta\ell/2}.
\]
All these statements, including the unique-ground-vector specialization,
are proved without polynomial growth. For $R=0$, the separate zero-range
theorem proves that both filters fix $H_i$ exactly and the ground-space
leakage is zero.

The same adjustment does not establish the uniform basis sum. The spectral
filtering rate depends on the gap and velocity and can be smaller than the
exponential volume-growth rate; repeated spectral splits weaken it further.
The preceding matrix family excludes an argument based only on the present
pointwise decay estimate, and the physical construction cited above shows
that the general metric version of eq.~(oibu) is false. The polynomial-growth
basis theorem is the proved corrected statement. The resolved verdict records
two separate outcomes: eq.~(facld) now holds with the source's general metric
hypotheses, while the unrestricted basis assertion is false and its
polynomial-growth correction is proved.

\bibliographystyle{unsrt}
\bibliography{references}
\end{document}
